\documentclass[12pt]{article}
\usepackage[utf8]{inputenc}
\usepackage[T1]{fontenc}
\usepackage[margin=1in]{geometry}
\usepackage{setspace}
\usepackage{xcolor}
\usepackage[colorlinks=true, urlcolor=blue!60!black, linkcolor=blue!60!black]{hyperref}
\setlength{\parindent}{0em}
\setlength{\parskip}{0.6em}
\renewcommand{\baselinestretch}{1.0}
\begin{document}
\begin{center}
{\LARGE \textit{Employment Division v. Smith}, 494 U.S. 872 (1990)\par}
\vspace{0.5em}
{\large Case Brief --- Two-Page Edition\par}
\end{center}

\section{Memory Jogger}
The Free Exercise Clause does not exempt individuals from valid, neutral laws of general applicability, even when those laws incidentally burden central religious practices such as the sacramental use of peyote (\href{https://supreme.justia.com/cases/federal/us/494/872/}{494 U.S. 872}).

\section{Detailed Case Facts}
Alfred Smith and Galen Black, drug-rehabilitation counselors at a private Oregon nonprofit, were fired for ingesting peyote at a Native American Church ceremony; their sincerity was never disputed. Oregon denied them unemployment benefits for work-related ``misconduct.'' Oregon law made possessing peyote, a Schedule I substance, a felony, though Oregon never prosecuted them.

\section{Procedural History}
The Oregon courts held the denial unconstitutional under \textit{Sherbert v. Verner}. In \textit{Smith I} (1988) the Supreme Court remanded to determine whether Oregon criminal law banned sacramental peyote use. The Oregon Supreme Court held that it did but that the ban violated the Free Exercise Clause. The Supreme Court reversed.

\section{Judicial Votes}
Reversed, 6--3, but 5--4 on the standard. Scalia wrote for Rehnquist, White, Stevens, and Kennedy; O'Connor concurred in the judgment, with Brennan, Marshall, and Blackmun joining her defense of strict scrutiny; Blackmun, with Brennan and Marshall, dissented.

\section{Holding}
The Free Exercise Clause does not relieve an individual of the obligation to comply with a ``valid and neutral law of general applicability.'' \textit{Sherbert}'s compelling-interest test does not apply to such laws, and religious exemptions are left to the political process. The ``hybrid rights'' and ``individual exemptions'' carve-outs were arguably not necessary to the decision.

\section{Analysis of Opinions}
Scalia reasoned that an incidental burden from a neutral law is not prohibiting'' free exercise; that the Court had never exempted religious conduct from such a law on free exercise grounds alone (recasting \textit{Yoder} and \textit{Cantwell} as hybrid cases); that \textit{Sherbert} applied only to individualized assessments like unemployment good cause''; and that strict scrutiny for every objection would make each conscience a law unto itself. O'Connor would have kept strict scrutiny but found the peyote ban survived it. Blackmun found Oregon's interest merely symbolic, noting no prosecutions, no evidence of harm, negligible trafficking, and federal and state exemptions.

\section{Examples: Future Applications}
\textbf{Same-Side:} a uniform school vaccination mandate with only medical exemptions requires no religious exemption. \textbf{Opposite-Side:} a foster-care contract allowing discretionary exceptions triggers strict scrutiny when an exception is denied to a religious agency, as in \textit{Fulton}. \textbf{Fence-Sitter:} a religious web designer raises a hybrid'' free exercise and compelled-speech claim against a neutral public-accommodations law.

\section{Critique}
\textit{Smith} provoked a bipartisan backlash that produced RFRA, then RLUIPA and state RFRAs. Scholars call the hybrid-rights distinction unprincipled and the individualized-exemptions carve-out unstable, and originalists such as Michael McConnell argue the Founders understood free exercise to require exemptions, while Philip Hamburger and others defend the result.

\section{Key Quotations}
\begin{quote}
``We have never held that an individual's religious beliefs excuse him from compliance with an otherwise valid law prohibiting conduct that the State is free to regulate.''
\end{quote}
\noindent --- Justice Scalia, Opinion of the Court, 494 U.S. at 878--879.

\end{document}
